\scanpage{82}
est un isomorphisme \inserted{en cohomologie} en bidegré $(i,j)$ avec $j\leq i$,
tandis que la coh. du premier mb est nulle pour $j>i$.

On peut l'exprimer comme un isomorphisme du $1^{\mathrm{er}}$ mb avec un certain,
tronqué du second $\ldots$

\uline{Fonctorialités}\quad a) en $k$, clair

b) \uline{En $u$}\deleted{\unsure{k}} Soit
\[
u':\mathbb Z\longrightarrow\Phi'
\]
une deuxième immersion; S'il existe un hom $\lambda:\Phi\longrightarrow\Phi'$
rendant comm.
\diagram{\begin{tikzcd}[column sep=large]
&\mathbb Z\arrow[dl,"u"']\arrow[dr,"u'"]&\\
\Phi\arrow[rr,"\lambda"]&&\Phi'
\end{tikzcd}}
i.e. $\lambda u=u'$

On en conclut $\unsure{\bar\lambda|\tilde\lambda}:K_u^{\bullet\bullet}-K_{u'}^{\bullet\bullet}$
et $\tilde\lambda_k:K_{u,k}^{\bullet\bullet}-K_{u',k}^{\bullet\bullet}$,
compatible avec ttes les structures, qui est un isom. dans \unsure{telle|toute}
catégorie dérivée, et le reste après application de $\Gamma(X,-)$ ou de $f_*$.

Supposons qu'on ait un autre, soit $\mu$, tel que
\[
\mu u=u'
\]
alors, posant $\mu=\lambda+\unsure{\theta|\Theta}$,
on trouve
\[
\unsure{\theta|\Theta}u=0
\]
\margin{81}
